% ============================================================
% CAPÍTULO 2: ESTADO DEL ARTE
% ============================================================

\section{El Método Socrático}
% Describe qué es el método socrático, su origen y aplicación en la educación.
% - Diálogo cooperativo basado en preguntas
% - Pensamiento crítico
% - Aplicaciones históricas y modernas
% TODO: Redactar


\section{Inteligencia Artificial en la Educación}
% Analiza las herramientas de IA existentes usadas en educación:
% - GitHub Copilot, ChatGPT, Gemini, etc.
% - Ventajas e inconvenientes para el aprendizaje
% - Comparativa con tu propuesta
% TODO: Redactar


\section{Modelos de Lenguaje Extenso (LLMs)}
% Explica qué son los LLMs, cómo funcionan a alto nivel, y cuáles son
% los modelos relevantes para tu proyecto (OpenAI, Anthropic, modelos abiertos).
% TODO: Redactar
% Referencia inicial para que BibTeX procese la bibliografía:
Entre los modelos más relevantes destaca GPT-4~\cite{openai2023gpt4}.


\section{Ingeniería de \textit{Prompts} (Prompt Engineering)}
% Describe las técnicas de prompt engineering relevantes:
% - System prompts
% - Few-shot prompting
% - Restricciones de comportamiento
% TODO: Redactar


\section{Tecnologías Web Relevantes}
% Revisa brevemente las tecnologías usadas en el proyecto:
% - React / frameworks frontend
% - Node.js / Express / frameworks backend
% - Bases de datos (PostgreSQL, etc.)
% - APIs de LLMs
% TODO: Redactar
